% !TITLE 青青陵上柏
% !AUTHOR 无名氏
% !DYNASTY 两汉
% !GENRE 五言古诗
% !ORDER 2
% !SOURCE 《古诗十九首》；收于《文选》
% !EDITION 《古诗十九首集释》，隋树森，中华书局（2020年版为候选）
% !EDITION_STATUS pending

\begin{poem}
青青陵上柏\textsuperscript{1}，磊磊涧中石。
人生天地间，忽如远行客\textsuperscript{2}。
斗酒相娱乐，聊厚不为薄\textsuperscript{3}。
驱车策驽马\textsuperscript{4}，游戏宛与洛\textsuperscript{5}。
洛中何郁郁，冠带自相索\textsuperscript{6}。
长衢罗夹巷，王侯多第宅。
两宫遥相望，双阙百余尺\textsuperscript{7}。
极宴娱心意，戚戚何所迫\textsuperscript{8}？
\end{poem}

\begin{annotations}
\notetext{1}{陵：大土山。磊磊：乱石堆积的样子。涧：山间的水沟。}
\notetext{2}{忽：迅速、飘忽。远行客：远行的旅人。人活在天地之间，像住店的旅客一样短暂。}
\notetext{3}{聊厚不为薄：姑且把薄酒当作厚酒，不嫌它寡淡。聊，姑且。厚、薄，指酒的醇厚与淡薄。}
\notetext{4}{策：鞭打、驱赶。驽马：劣马。}
\notetext{5}{宛：宛县，东汉南阳郡治所，在今河南南阳。洛：洛阳，东汉都城。}
\notetext{6}{郁郁：繁盛热闹的样子。冠带：戴冠束带的人，指士大夫、官吏。自相索：互相探访、结交。索，寻求。}
\notetext{7}{长衢：宽阔的大道。夹巷：大道两旁的小巷。第宅：贵族的宅院。两宫：洛阳城内的南宫与北宫。双阙：宫殿门前两侧的高台建筑。}
\notetext{8}{极宴：尽情宴饮。戚戚：忧愁不安的样子。何所迫：被什么逼迫着。}
\end{annotations}

\begin{appreciation}
《青青陵上柏》是古诗十九首的第三首，写一个漂泊者在洛阳的所见所感。

诗从两样东西写起：“青青陵上柏，磊磊涧中石”——山陵上柏树青青，山涧里乱石磊磊。柏树耐寒，石头不朽，它们都比人活得长久。人站在这两样东西中间，忽然就觉出了自己的渺小：“人生天地间，忽如远行客。”人在天地之间，就像一位匆匆赶路的旅客，住店一夜，天亮就走。这个比喻看似平静，底下却是极大的惊动——原来我们不是这里的主人，只是过客。

既然只是过客，那就喝酒吧。“斗酒相娱乐，聊厚不为薄”——一斗薄酒也可以自得其乐，姑且把它当作厚酒，不嫌它薄。这是穷人的自我宽解，也是乱世里的处世之道。喝完酒，诗人驱车出游：“驱车策驽马，游戏宛与洛。”驾着劣马，在宛县和洛阳之间游荡。驽马二字很实在，他连一匹好马都没有。

接下来六句是全诗最热闹的地方。洛阳城里何等繁盛：“洛中何郁郁，冠带自相索”——满城冠带，官员士大夫们互相拜访、彼此攀结。大道两旁罗列着小巷，王侯的宅第鳞次栉比；南宫北宫遥遥相望，宫门前的双阙高达百余尺。诗人用极省净的笔，把东汉都城的气派一层层铺开：从街市到宅第，从宅第到宫阙，视线一路向上，最后停在“百余尺”的高度上。

可是就在这极宴之中，他忽然问了一句：“极宴娱心意，戚戚何所迫？”——尽情宴饮本该使人心意畅快，为什么我心里还是戚戚不安，像被什么东西逼迫着？这一问把前面所有的繁华都翻了过来。满城的富贵是人家的，他只有一斗薄酒、一匹驽马。热闹越是铺得足，这一问就越是冷。

同样是写洛阳，班固的《两都赋》、张衡的《二京赋》用的是铺天盖地的繁辞，宫室、苑囿、物产、礼仪，一样一样地数；这首诗只用了六句，却让读者记住了那些宅第和双阙，也记住了站在它们外面的那个人。赋写的是城，诗写的是人。

结尾那一问没有答案。两千年前的洛阳早已埋在土里，柏树还在长，涧石还在，过客换了一批又一批。
\end{appreciation}
